\begin{exercisechunk}
\item \textbf{Sampling with and without replacement.}
\label[exercise]{ex:l6-resampling}
\exerciseprofile{2}{2}{3}{\exproof\quad\excalculation\quad\exinterpretation}{%
Derive the laws induced by sampling indices with and without replacement and
identify what these procedures do not change about the source data.}
Fix a corpus \(x_1,\ldots,x_m\in\R\) and let
\(\bar x=m^{-1}\sum_i x_i\).
\begin{enumerate}[(a),beginpenalty=10000]
\item Draw \(I_1,\ldots,I_B\) independently and uniformly from \([m]\), and
define \(C_i=\sum_{b=1}^B\I\{I_b=i\}\). Give the joint law of
\((C_1,\ldots,C_m)\), and prove that
\(B^{-1}\sum_iC_ix_i\) is conditionally unbiased for \(\bar x\).
\item Take \(B=m\). Find the marginal law of \(C_i\) and its limit as
\(m\to\infty\). Find the limiting probability that item \(i\) is omitted.
Why does this calculation not show that the resampled items are independent
population observations?
\item Suppose instead that indices are sampled without replacement. Prove that
sampling all \(m\) indices gives a uniformly random permutation. Show that the
first \(B<m\) indices are a uniformly random ordered sample without replacement,
or equivalently a uniformly random \(B\)-element subset followed by a uniformly
random ordering.
\end{enumerate}
\end{exercisechunk}
